Generative audio systems can now synthesize speech, singing, music, and environmental sound with increasing realism, which threatens media integrity and trust in voice-based communication. A practical detector must work when the audio category is unknown at test time and when recordings have passed through lossy compression, resampling, noise, or narrowband channels that alter the spectral cues classifiers rely on. It must also be cheap to run: large self-supervised encoders achieve strong benchmark scores but are difficult to deploy on mobile or embedded hardware, which motivates compact models whose robustness is improved during training rather than by enlarging the inference graph.

This report investigates lightweight audio deepfake detection on the AT-ADD Track~2 benchmark. A MobileNetV3-Small baseline (${\sim}1.52$\,M parameters) operating on 80-bin log-mel features at 24\,kHz is trained on clean audio only and compared with \textbf{Frequency-Time Robustness Coherence Learning (FT-RCL)}, a training framework that exposes the same backbone to paired clean and synthetically degraded waveforms. FT-RCL combines a dual classification loss, prediction coherence, eight-band feature signature coherence, a severity curriculum, and gated robustness weights, and it leaves the inference graph unchanged. A STAR-Edge extension, which adds asymmetric bona-fide robustness, severity trajectory coherence, and channel auxiliary detection, is also implemented. On 13{,}287 clean held-out clips, the baseline attains a Macro-F1 of 0.672 and an equal error rate (EER) of 0.216, compared with 0.659 and 0.234 for FT-RCL. FT-RCL improves Macro-F1 on environmental sound and music but yields no aggregate advantage on clean audio. A preliminary conditioned evaluation on 100 clips per model reports $\Delta$Macro-F1 under resampling, noise, and telephony band-limiting; full-split $\Delta M$ analysis remains future work.
